\documentclass[border=12pt]{standalone}
\usepackage[no-math]{fontspec}
\usepackage[brazil]{babel}
\usepackage{amsmath,amssymb}
\IfFileExists{C:/Windows/Fonts/Lato_Regular.ttf}{%
  \setmainfont{Lato}[
    Path = C:/Windows/Fonts/,
    UprightFont = Lato_Regular.ttf,
    BoldFont = Lato_Bold.ttf,
    ItalicFont = Lato_Italic.ttf,
    BoldItalicFont = Lato_BoldItalic.ttf
  ]
  \setsansfont{Lato}[
    Path = C:/Windows/Fonts/,
    UprightFont = Lato_Regular.ttf,
    BoldFont = Lato_Bold.ttf,
    ItalicFont = Lato_Italic.ttf,
    BoldItalicFont = Lato_BoldItalic.ttf
  ]
}{%
  \RequirePackage[default]{lato}
}
\renewcommand{\familydefault}{\sfdefault}

\usepackage{xcolor}
\usepackage{tikz}
\usepackage{pgfplots}
\pgfplotsset{compat=1.18}
\usetikzlibrary{arrows.meta,positioning,shapes.geometric,fit,backgrounds,calc,decorations.pathreplacing,patterns}

% Paleta institucional UEFS / Extensão Rural
\definecolor{erNavy}{HTML}{1F4268}
\definecolor{erNavyDark}{HTML}{152E4A}
\definecolor{erGold}{HTML}{B38A3D}
\definecolor{erGoldDark}{HTML}{8C6927}
\definecolor{erTeal}{HTML}{0E7C7B}
\definecolor{erForest}{HTML}{2D6A4F}
\definecolor{erAlert}{HTML}{9E2A2B}
\definecolor{erInk}{HTML}{26313D}
\definecolor{erMuted}{HTML}{5B6573}

\definecolor{bgNavy}{HTML}{EDF2F7}
\definecolor{bgGold}{HTML}{F8F3E9}
\definecolor{bgTeal}{HTML}{EBF6F5}
\definecolor{bgForest}{HTML}{EEF5F0}
\definecolor{bgAlert}{HTML}{FDF2F2}

\tikzset{
  >=Stealth,
  erBox/.style={
    draw=#1, fill=white, rounded corners=5pt, line width=1.15pt,
    inner sep=7pt, align=center, font=\footnotesize\sffamily, text=erInk
  },
  erTint/.style={
    draw=#1, fill=#1!8!white, rounded corners=5pt, line width=1.15pt,
    inner sep=7pt, align=center, font=\footnotesize\sffamily, text=erInk
  },
  erBadge/.style={
    fill=#1, text=white, rounded corners=9pt, inner xsep=8pt, inner ysep=3pt,
    font=\scriptsize\bfseries\sffamily
  },
  erTitle/.style={font=\normalsize\bfseries\sffamily, text=erNavy, align=center},
  erSub/.style={font=\footnotesize\sffamily, text=erMuted, align=center},
  erArrow/.style={->, >=Stealth, line width=1.35pt},
  erDArrow/.style={<->, >=Stealth, line width=1.35pt}
}

\hyphenpenalty=10000
\exhyphenpenalty=10000

\begin{document}
\begin{tikzpicture}[font=\sffamily]

% Título principal
\node[erTitle, font=\large\bfseries\sffamily] (title) at (0, 5.15)
  {Avaliação Participativa de Impacto em Tecnologias Sociais};

% Subtítulo da faixa temporal
\node[erSub, font=\scriptsize\bfseries\sffamily, text=erMuted] (tl_label) at (0, 4.45)
  {CICLO TEMPORAL DA AVALIAÇÃO};

% Fluxo temporal superior
\node[
  erTint=erNavy, text width=4.35cm, minimum height=1.25cm, inner sep=8pt
] (t1) at (-5.6, 3.35) {%
  {\color{erNavy}\small\bfseries Ex-ante}\\[3pt]
  {\color{erInk}\footnotesize Linha de Base}
};

\node[
  erTint=erTeal, text width=4.65cm, minimum height=1.25cm, inner sep=8pt
] (t2) at (0, 3.35) {%
  {\color{erTeal}\small\bfseries De Processo}\\[3pt]
  {\color{erInk}\footnotesize Monitoramento Participativo}
};

\node[
  erTint=erForest, text width=4.35cm, minimum height=1.25cm, inner sep=8pt
] (t3) at (5.6, 3.35) {%
  {\color{erForest}\small\bfseries Ex-post}\\[3pt]
  {\color{erInk}\footnotesize Impacto Territorial}
};

\draw[erArrow, color=erNavy!80] (t1.east) -- (t2.west);
\draw[erArrow, color=erTeal!80] (t2.east) -- (t3.west);

% Linha divisória sutil para comparação
\draw[color=erMuted!35, line width=0.9pt] (-8.0, 2.25) -- (8.0, 2.25);
\node[
  fill=white, text=erMuted, font=\scriptsize\bfseries\sffamily, inner xsep=8pt
] at (0, 2.25) {CONTRASTE DE ABORDAGENS};

% Card esquerdo: Avaliação Convencional (Top-Down)
\node[
  draw=erAlert, fill=bgAlert, rounded corners=6pt, line width=1.25pt,
  text width=6.55cm, inner sep=11pt, align=left, anchor=north
] (conv) at (-4.3, 1.75) {%
  \parbox[t][3.25cm][t]{6.55cm}{%
    {\color{erAlert}\small\bfseries Avaliação Convencional (Top-Down)}\\[4pt]
    {\color{erAlert!45}\rule{6.55cm}{0.8pt}}\\[8pt]
    {\color{erInk}\footnotesize
    \textbf{\color{erAlert}$\times$}\enspace Ignora o saber local\\[6pt]
    \textbf{\color{erAlert}$\times$}\enspace Mede apenas produtividade e renda\\[6pt]
    \textbf{\color{erAlert}$\times$}\enspace Não capta autonomia ou gênero\\[6pt]
    \textbf{\color{erAlert}$\times$}\enspace Sem retorno à comunidade}
  }%
};

% Card direito: Avaliação Participativa (Co-investigação)
\node[
  draw=erForest, fill=bgForest, rounded corners=6pt, line width=1.25pt,
  text width=6.55cm, inner sep=11pt, align=left, anchor=north
] (part) at (4.3, 1.75) {%
  \parbox[t][3.25cm][t]{6.55cm}{%
    {\color{erForest}\small\bfseries Avaliação Participativa (Co-investigação)}\\[4pt]
    {\color{erForest!45}\rule{6.55cm}{0.8pt}}\\[8pt]
    {\color{erInk}\footnotesize
    \textbf{\color{erForest}\checkmark}\enspace Beneficiários atuam como co-avaliadores\\[6pt]
    \textbf{\color{erForest}\checkmark}\enspace Indicadores quantitativos e qualitativos\\[6pt]
    \textbf{\color{erForest}\checkmark}\enspace Empoderamento e controle social}
  }%
};

% Selo central de contraste
\node[
  circle, fill=white, draw=erMuted!60, line width=1pt,
  font=\scriptsize\bfseries\sffamily, text=erMuted, inner sep=4.5pt
] at (0, -0.25) {vs.};

\end{tikzpicture}
\end{document}
